\chapter{2022年北京卷}

\section{单选题}

\begin{problem}
已知全集 \(U=\{x\mid -3<x<3\}\),集合 \(A=\{x\mid -2<x\le 1\}\),则 \(C_UA=\)(\quad).

\choices
    {\((-2,1]\)}
    {\((-3,-2)\cup[1,3)\)}
    {\([-2,1)\)}
    {\((-3,-2]\cup(1,3)\)}
\end{problem}

\begin{problem}
若复数 \(z\) 满足 \(\i z=3-4\i\),则 \(|z|=\)(\quad).

\choices
    {1}
    {5}
    {7}
    {25}
\end{problem}

\begin{problem}
若直线 \(2x+y-1=0\) 是圆 \((x-a)^2+y^2=1\) 的一条对称轴,则 \(a=\)(\quad).

\choices
    {\(\frac12\)}
    {\(-\frac12\)}
    {1}
    {\(-1\)}
\end{problem}

\begin{problem}
已知函数 \(f(x)=\dfrac{1}{1+2^x}\),则对任意实数 \(x\),有(\quad).

\choices
    {\(f(-x)+f(x)=0\)}
    {\(f(-x)-f(x)=0\)}
    {\(f(-x)+f(x)=1\)}
    {\(f(-x)-f(x)=\frac13\)}
\end{problem}

\begin{problem}
已知函数 \(f(x)=\cos^2x-\sin^2x\),则(\quad).

\choices
    {\(f(x)\) 在 \(\left(-\frac{\pi}{2},-\frac{\pi}{6}\right)\) 上单调递减}
    {\(f(x)\) 在 \(\left(-\frac{\pi}{4},\frac{\pi}{12}\right)\) 上单调递增}
    {\(f(x)\) 在 \(\left(0,\frac{\pi}{3}\right)\) 上单调递减}
    {\(f(x)\) 在 \(\left(\frac{\pi}{4},\frac{7\pi}{12}\right)\) 上单调递增}
\end{problem}

\begin{problem}
设 \(\{a_n\}\) 是公差不为 0 的无穷等差数列,则``\(\{a_n\}\) 为递增数列''是``存在正整数 \(N_0\),当 \(n>N_0\) 时,\(a_n>0\)''的(\quad).

\choices
    {充分而不必要条件}
    {必要而不充分条件}
    {充分必要条件}
    {既不充分也不必要条件}
\end{problem}

\begin{problem}
在北京冬奥会上,国家速滑馆``冰丝带''使用高效环保的二氧化碳跨临界直冷制冰技术,为实现绿色冬奥作出了贡献. 如图描述了一定条件下二氧化碳所处的状态与 \(T\) 和 \(\lg P\) 的关系,其中 \(T\) 表示温度,单位是 \(\mathrm{K}\);\(P\) 表示压强,单位是 \(\mathrm{bar}\). 下列结论中正确的是(\quad).

\begin{center}
\bitmapfigure[width=0.34\textwidth]{img/2022/beijing/q7_fig1.png}
\end{center}

\choices
    {当 \(T=220\),\(P=1026\) 时,二氧化碳处于液态}
    {当 \(T=270\),\(P=128\) 时,二氧化碳处于气态}
    {当 \(T=300\),\(P=9987\) 时,二氧化碳处于超临界状态}
    {当 \(T=360\),\(P=729\) 时,二氧化碳处于超临界状态}
\end{problem}

\begin{problem}
若 \((2x-1)^4=a_4x^4+a_3x^3+a_2x^2+a_1x+a_0\),则 \(a_0+a_2+a_4=\)(\quad).

\choices
    {40}
    {41}
    {\(-40\)}
    {\(-41\)}
\end{problem}

\begin{problem}
已知正三棱锥 \(P-ABC\) 的六条棱长均为 6,\(S\) 是 \(\triangle ABC\) 及其内部的点构成的集合. 设集合 \(T=\{Q\in S\mid PQ\le 5\}\),则 \(T\) 表示的区域的面积为(\quad).

\choices
    {\(\frac{3\pi}{4}\)}
    {\(\pi\)}
    {\(2\pi\)}
    {\(3\pi\)}
\end{problem}

\begin{problem}
在 \(\triangle ABC\) 中,\(AC=3\),\(BC=4\),\(\angle C=90^\circ\). \(P\) 为 \(\triangle ABC\) 所在平面内的动点,且 \(PC=1\),则 \(\overrightarrow{PA}\cdot \overrightarrow{PB}\) 的取值范围是(\quad).

\choices
    {\([-5,3]\)}
    {\([-3,5]\)}
    {\([-6,4]\)}
    {\([-4,6]\)}
\end{problem}

\section{填空题}

\begin{problem}
函数 \(f(x)=\dfrac{1}{x}+\sqrt{1-x}\) 的定义域是\fillinblank{}.
\end{problem}

\begin{problem}
已知双曲线 \(y^2+\dfrac{x^2}{m}=1\) 的渐近线方程为 \(y=\pm \dfrac{\sqrt3}{3}x\),则 \(m=\)\fillinblank{}.
\end{problem}

\begin{problem}
若函数 \(f(x)=A\sin x-\sqrt3\cos x\) 的一个零点为 \(\dfrac{\pi}{3}\),则 \(A=\)\fillinblank{};\(f\left(\dfrac{\pi}{12}\right)=\)\fillinblank{}.
\end{problem}

\begin{problem}
设函数 \(f(x)=\begin{cases}
-ax+1,&x<a,\\
\left(x-2\right)^2,&x\ge a
\end{cases}\),若 \(f(x)\) 存在最小值,则 \(a\) 的一个取值为\fillinblank{};\(a\) 的最大值为\fillinblank{}.
\end{problem}

\begin{problem}
已知数列 \(\{a_n\}\) 各项均为正数,其前 \(n\) 项和 \(S_n\) 满足 \(a_nS_n=9\)(\(n=1\),\(2\),\(\dots\)). 给出下列四个结论:

\begin{circlelist}
\item \(\{a_n\}\) 的第 2 项小于 3;

\item \(\{a_n\}\) 为等比数列;

\item \(\{a_n\}\) 为递减数列;

\item \(\{a_n\}\) 中存在小于 \(\dfrac{1}{100}\) 的项.
\end{circlelist}

其中所有正确结论的序号是\fillinblank{}.
\end{problem}

\section{解答题}

\begin{problem}
在 \(\triangle ABC\) 中,\(\sin 2C=\sqrt3\sin C\).

\begin{enumerate}
\item 求 \(\angle C\);

\item 若 \(b=6\),且 \(\triangle ABC\) 的面积为 \(6\sqrt3\),求 \(\triangle ABC\) 的周长.
\end{enumerate}
\end{problem}

\begin{problem}
如图,在三棱柱 \(ABC-A_1B_1C_1\) 中,侧面 \(BCC_1B_1\) 为正方形,平面 \(BCC_1B_1\perp\) 平面 \(ABB_1A_1\),\(AB=BC=2\),\(M\),\(N\) 分别为 \(A_1B_1\),\(AC\) 的中点.



\begin{enumerate}
\item 求证:\(MN\parallel\) 平面 \(BCC_1B_1\);

\item 再从下列两个条件中选择一个作为已知,求直线 \(AB\) 与平面 \(BMN\) 所成角的正弦值.

\begin{circlelist}
\item \(AB\perp MN\);

\item \(BM=MN\).
\end{circlelist}

\end{enumerate}

\bitmapfigure[width=0.30\textwidth]{img/2022/beijing/q17_fig1.png}
\end{problem}


\begin{problem}
在校运动会上,只有甲,乙,丙三名同学参加铅球比赛,比赛成绩达到 \(9.50\text{ m}\) 以上(含 \(9.50\text{ m}\))的同学将获得优秀奖. 为预测获得优秀奖的人数及冠军得主,收集了甲,乙,丙以往的比赛成绩,并整理得到如下数据(单位:\(\mathrm{m}\)):

甲:\(9.80\),\(9.70\),\(9.55\),\(9.54\),\(9.48\),\(9.42\),\(9.40\),\(9.35\),\(9.30\),\(9.25\);

乙:\(9.78\),\(9.56\),\(9.51\),\(9.36\),\(9.32\),\(9.23\);

丙:\(9.85\),\(9.65\),\(9.20\),\(9.16\).

假设用频率估计概率,且甲,乙,丙的比赛成绩相互独立.

\begin{enumerate}
\item 估计甲在校运动会铅球比赛中获得优秀奖的概率;

\item 设 \(X\) 是甲,乙,丙在校运动会铅球比赛中获得优秀奖的总人数,估计 \(X\) 的数学期望 \(E(X)\);

\item 在校运动会铅球比赛中,甲,乙,丙谁获得冠军的概率估计值最大?(结论不要求证明).
\end{enumerate}
\end{problem}


\begin{problem}
已知椭圆 \(E:\frac{x^2}{a^2}+\frac{y^2}{b^2}=1\)(\(a>b>0\))的一个顶点为 \(A(0,1)\),焦距为 \(2\sqrt3\).

\begin{enumerate}
\item 求椭圆 \(E\) 的方程;

\item 过点 \(P(-2,1)\) 作斜率为 \(k\) 的直线与椭圆 \(E\) 交于不同的两点 \(B\),\(C\),直线 \(AB\),\(AC\) 分别与 \(x\) 轴交于点 \(M\),\(N\). 当 \(|MN|=2\) 时,求 \(k\) 的值.
\end{enumerate}
\end{problem}


\begin{problem}
已知函数 \(f(x)=\e^x\ln(1+x)\).

\begin{enumerate}
\item 求曲线 \(y=f(x)\) 在点 \((0,f(0))\) 处的切线方程;

\item 设 \(g(x)=f'(x)\),讨论函数 \(g(x)\) 在 \([0,+\infty)\) 上的单调性;

\item 证明:对任意的 \(s\),\(t\in(0,+\infty)\),有 \(f(s+t)>f(s)+f(t)\).
\end{enumerate}
\end{problem}


\begin{problem}
已知 \(Q:a_1\),\(a_2\),\(\dots\),\(a_k\) 为有穷整数数列. 给定正整数 \(m\),若对任意的 \(n\in\{1,2,\dots,m\}\),在 \(Q\) 中存在 \(a_i\),\(a_{i+1}\),\(\dots\),\(a_{i+j}\)(\(j\ge 0\)),使得 \(a_i+a_{i+1}+\dots+a_{i+j}=n\),则称 \(Q\) 为 \(m\)-连续可表数列.

\begin{enumerate}
\item 判断 \(Q:2\),\(1\),\(4\) 是否为 5-连续可表数列?是否为 6-连续可表数列?说明理由;

\item 若 \(Q:a_1\),\(a_2\),\(\dots\),\(a_k\) 为 8-连续可表数列,求证:\(k\) 的最小值为 4;

\item 若 \(Q:a_1\),\(a_2\),\(\dots\),\(a_k\) 为 20-连续可表数列,且 \(a_1+a_2+\dots+a_k<20\),求证:\(k\ge 7\).
\end{enumerate}
\end{problem}

